\chapter{Conclusions}



\section{Further improvements} \label{lab:c6-further-improvements}

\app{} does not support TPM operations on Windows.
While I would have loved to make it work on Windows as well, it simply looks like TPM development is not meant to be done on Windows unless you are a TPM expert.
The \t{tpm2-pytss} library does not support Windows\footnote{\url{https://github.com/tpm2-software/tpm2-pytss/issues/597\#issuecomment-2556916125}}, and there are no similar libraries for Windows.
So writing equivalent FAPI code was not an option.
The remaining solution was to go low level, calling the TPM 2.0 commands directly (\t{TPM2\_NV\_ReadPublic}, \t{TPM2\_GetRandom}, \t{TPM2\_NV\_DefineSpace}, \t{TPM2\_NV\_Write}, \t{TPM2\_NV\_Read}).
The only way I found to call these was through the \t{Tbsip\_Submit\_Command} C API.\footnote{\url{https://learn.microsoft.com/en-us/windows/win32/api/tbs/nf-tbs-tbsip_submit_command}}
This API expects you to build the commands byte by byte.
I found no examples online of how to use it.
AI sketched me an example of how the code working on Linux would look like with calls to this API, and it gave me more than 200 lines of byte-manipulating code.
Unfortunately, my current TPM knowledge is not deep enough to write and document such code, nor do I wish to write such low-level code for something that should be a few lines, but isn't, since Windows.
I also had the idea to use WSL, but it looks like WSL has no access to the physical TPM.
